\section{Matched block maps for product algebras}

The Fundamental Theorem compares the simple matrix summands selected by a
product-algebra equivalence.  The following tensor-network-facing statements
record the component maps and their dimension consequence.

\begin{theorem}[Paired component maps for products of matrix algebras]
    \label{thm:mps_component_map_bijective}
    \lean{TwoSidedIdeal.ringEquiv_maps_single_support_between}
    \lean{TwoSidedIdeal.blockComponentMap_surjective}
    \lean{TwoSidedIdeal.blockComponentMap_bijective}
    \leanok

    For a ring equivalence between products of full matrix algebras, the map
    induced between each pair of matched components is bijective.
\end{theorem}

\begin{proof}\leanok

    Apply Theorem~\cite[``Paired component maps are bijections'']{QICLean2026Blueprint} to the matched block
    ideals.
\end{proof}

\begin{theorem}[Dimension preservation for matched matrix summands]
    \label{thm:mps_dim_preserved}
    \lean{MPSTensor.dim_preserved}
    \leanok
    An algebra automorphism of a product of full matrix algebras preserves
    the dimension of every pair of summands matched by its induced
    permutation.
\end{theorem}

\begin{proof}\leanok
    \uses{thm:mps_component_map_bijective}
    Restrict the automorphism to the two matched summands. The resulting
    component map is a bijective complex-linear map by
    Theorem~\ref{thm:mps_component_map_bijective}. Hence the two matrix
    spaces have equal complex dimensions, so their sizes have equal squares
    and therefore are equal.
\end{proof}
